% Figure: a force in three dimensions, resolved into Cartesian components
% by its direction cosines. The angles the force makes with each axis fix
% the unit vector; the cosines square-sum to one.
\begin{figure}[htbp]
\centering
\begin{tikzpicture}[
  axis/.style={draw=engblack, -{Stealth}, thick},
  force/.style={draw=engblue, -{Stealth}, very thick},
  comp/.style={draw=enggray, thin, dashed},
  scale=1.0,
]
% an oblique 3D frame: x toward lower-right, y toward right, z up
\coordinate (O) at (0,0);
\coordinate (Xax) at (-2.3,-1.5);   % x axis (out of page, down-left)
\coordinate (Yax) at (4.0,0);       % y axis (to the right)
\coordinate (Zax) at (0,3.6);       % z axis (up)
\draw[axis] (O) -- (Xax) node[anchor=north east, font=\small] {$x$};
\draw[axis] (O) -- (Yax) node[anchor=west, font=\small] {$y$};
\draw[axis] (O) -- (Zax) node[anchor=south, font=\small] {$z$};
% the force tip P, an oblique point in the box
\coordinate (P) at (1.9,2.5);
% projection of P onto the x-y "floor": move down from P by the z part
\coordinate (Q) at (1.9,0.55);      % foot of vertical from P onto floor
% the force vector
\draw[force] (O) -- (P) node[midway, above left, font=\small, engblue] {$\vec{F}$};
% vertical drop to the floor (z component)
\draw[comp] (P) -- (Q);
% floor projection from O to Q
\draw[comp] (O) -- (Q);
% small marks for the three axis angles
\draw[draw=engred, thin] (0.85,0) arc (0:53:0.85);
\node[font=\scriptsize, engred] at (1.15,0.45) {$\theta_z$};
% label the three direction angles schematically near O
\node[font=\scriptsize, anchor=west] at (1.6,1.45)
  {$\hat{u}=(\cos\theta_x,\cos\theta_y,\cos\theta_z)$};
% mark the foot
\fill[engblack] (Q) circle (0.04);
\node[font=\scriptsize, anchor=north west] at (Q) {$Q$};
\fill[engblue] (P) circle (0.04);
\node[font=\scriptsize, anchor=south west] at (P) {$P$};
\end{tikzpicture}
\caption[A force in three dimensions and its direction cosines]{A force
$\vec{F}$ in space is fixed by its magnitude and a unit vector
$\hat{u}=(\cos\theta_x,\cos\theta_y,\cos\theta_z)$, whose entries are the
cosines of the angles the force makes with the three axes. The point $P$
is the head of the force; $Q$ is its projection onto the $xy$ floor, and
the dashed drop $PQ$ is the vertical ($z$) component. The three cosines
are not independent: because $\hat{u}$ is a unit vector, they satisfy
$\cos^2\theta_x+\cos^2\theta_y+\cos^2\theta_z=1$, so any two of the angles
fix the third up to a sign. Resolving a spatial force is the same act as
the planar resolution of figure~\ref{fig:vol03ch01:force-vector}, carried
out three times rather than once.}
\label{fig:vol03ch01:direction-cosines}
\end{figure}
